\section{Supplying the actual Type-I hypotheses along a family}
\label{sec:type-one-family}

\status{GI1--GI8 prove the actual numerical hypothesis and its
uniform family version, then connect TC6 to the source's global
time and exceptional-share mean. GII supplies the good-orbit
count and segment geometry; GJ supplies the logarithmic size
bounds and the canonical exceptional function. The arithmetic
family and its stated size assumptions remain supplied data.}

\paragraph{GI1--GI2: the scales and one uniform threshold.}
Fix $0<\varepsilon<1/52$ and the nonnegative coefficient vector
$K$. Write $k=D.R.k$, $n=|HH(D)|$, and use the source definitions
\[
\begin{gathered}
 a=xe^{x^\varepsilon},\quad
 X'=e^{x^{\varepsilon/3}},\quad
 b'=a/X',\quad \delta_1=x^{-\varepsilon/8},\quad
 Y=x^{13},\quad Y_b=x^\varepsilon/20,\\
 \lambda=K.\mathrm{La}\,N_1^2\log a,\quad
 N_*=2n\delta_1,\quad
 c=4L.\mathrm{Sc}=16k\Delta,\quad
 I=\frac{N_*x^\varepsilon\lambda b'}k,\\
 \rho=\frac{4N_*x^\varepsilon\lambda}k
          +\frac{K.\rho\,N_1\lambda}k.
\end{gathered}                                                    \tag{GI.1}
\]
For a candidate prime $p$ and an actual good representative,
let $n'=\operatorname{ord}(h_p)$, let $T$ be its selected
positions, and put
\[
\begin{gathered}
 N=|T|,\qquad N_+=40n'\delta_1,\qquad Q=p^{G'},\\
 D_f=\left\lceil K.\mathrm{DI}\,N_+Y^2
                    (N_1^5+k^3)\log_p a\right\rceil_+.
\end{gathered}                                                    \tag{GI.2}
\]
The order of an element is at most the group order, and the
orbit parametrization is injective before its period. Thus
$N\le n'\le n$. The literal family size assumptions give
\[
 x^{9/2}\le k\le x^6,\qquad n\le x^3,\qquad
 \log\Delta\le\tfrac12x^{\varepsilon/4},\qquad
 c\le\lambda b'/k.                                      \tag{GI.3}
\]
GJ2, at $\eta=\varepsilon/48$, gives
\[
 x^{47\varepsilon/48}\le\lambda\le x^{49\varepsilon/48},
 \qquad \log Q\le x^{23\varepsilon/48}\quad(p<X').        \tag{GI.4}
\]
GII2 gives a fixed $B=B(\varepsilon)>0$ with
$G'\le B\log x\,x^{\varepsilon/8}$ and $N_1\le B\log x$.
One exact choice is
\[
 B=40\bigl((h+3)10^{6r+1}+6\bigr)(6r+1),\quad
 h=2+8\varepsilon,\quad
 r=\lceil8(h+2)/\varepsilon\rceil_++1.
\]
Enlarge a threshold depending only on $\varepsilon,K$ so that
\[
\begin{gathered}
 x\ge\max\{e,7,320K.\mathrm{DI}+1,K.\mathrm{kI}\},\\
 \max\{3B,514\}\log x\le x^{\varepsilon/8},\qquad
 (2B+3)\log x\le x,\qquad
 128\log x\le x^{47\varepsilon/48},\\
 20K.\mathrm{kI}\le x^\varepsilon,\qquad
 39\log x\le x^{\varepsilon/3},\qquad
 K.\mathrm{kI}\,x^4\le x^{9/2},\\
 340\log x\le x^\varepsilon,\quad
 16x^{23\varepsilon/48}\le x^\varepsilon,\quad
 2x^{\varepsilon/4}\le x^\varepsilon,\quad
 4x^{\varepsilon/3}\le x^\varepsilon.
\end{gathered}                                                    \tag{GI.5}
\]
Each inequality follows from a strict gap between exponents;
for example, putting $t=\log x$ and using
$e^{ut}\ge u^2t^2/2$ proves $C\log x\le x^u$ eventually.
The threshold precedes every prime, orbit and family index.

The candidate condition $n'\ge x^{\varepsilon/4}$ now gives
$3G'\le n'$. Also, for every prime $q$,
\[
 L_1(q)=\lceil\log_q((kx^{26})^4)\rceil_+
       \le\lceil256\log x\rceil_+,
\]
because $\log q\ge1/2$. Taking the bounded natural supremum
in the definition of $L_2$ gives
\[
 L_2<512\log x+2\le514\log x
       \le x^{\varepsilon/8}\le U=\lceil x^{\varepsilon/8}\rceil_+.
                                                               \tag{GI.6}
\]
The empty defining set gives $L_2=0$. These are exactly the two
fields of the source's larger \texttt{NumMean} record needed
here. GII1--GII3 now give the actual segment geometry and
$N\le N_+$.

\paragraph{GI3: verify the small numerical fields.}
Assume $p\ge p_0=k^2x^{26}$, $p<X'$ and $\rho\le x^{-2}$.
The first thirteen fields of \texttt{AuxBlock8.HypI} are
\[
\begin{gathered}
 k^2Y^2\le p,\quad p<a/b',\quad \rho<1/6,
 \quad 2N_1+2\le k-1,\quad 8\log(NY)\le\lambda,\\
 Y_bN\lambda b'/k\le I,\quad 2I\le b',\quad
 K.\mathrm{kI}\le Y_b,\quad
 kY^2Y_bN\lambda\log a\le X',\\
 K.\mathrm{kI}\,N\lambda^2\le k,\quad
 2Y_bN\lambda<k,\quad c\le\lambda b'/k,\quad N\le N_+.
\end{gathered}                                                    \tag{GI.7}
\]
Here is a direct check. Equations (GI.3)--(GI.5) give
$1\le\lambda\le x$, $\lambda^2\le x$,
$x^\varepsilon\lambda\le x$, $N_1\le x$ and
$\log a\le2x$. The rank margin follows from
$2N_1+3\le(2B+3)\log x\le x\le k$; the logarithmic bound
follows from $NY\le x^{16}$. The actual good orbits have $N>0$.
Moreover $N\le40n\delta_1$ gives $Y_bN\le N_*x^\varepsilon$.
The nonnegative summands of $\rho\le1/7$ give
$2I\le b'/14$, which is stronger than required.
The remaining product estimates are
\[
\begin{aligned}
 kY^2Y_bN\lambda\log a&\le2x^{38}\le x^{39}\le X',\\
 K.\mathrm{kI}\,N\lambda^2&\le K.\mathrm{kI}\,x^4\le k,\\
 2Y_bN\lambda&\le x^4/10<k.
\end{aligned}                                                    \tag{GI.8}
\]
The prime, conductor and top-count clauses are precisely their
stated hypotheses. Thus no one of the thirteen fields is an
additional asymptotic assumption.

\paragraph{GI4: the full depth and exponential budget.}
The fourteenth field is
\[
 K.\mathrm{kI}\,k^3NQ^4[8(D_f+1)+(N+n')c]Y\le b'.       \tag{GI.9}
\]
Since $N_+\le40x^3$, $N_1^5+k^3\le2x^{18}$ and
$\log_p a\le4x$, its actual rounded depth satisfies
\[
 D_f\le320K.\mathrm{DI}\,x^{48}+1\le x^{49}.
\]
Set $E_1=e^{x^{\varepsilon/4}/2}$ and
$E_2=e^{4x^{23\varepsilon/48}}$. Then $c\le16x^6E_1$,
$Q^4\le E_2$, and, for $x\ge7$,
\[
 8(D_f+1)+(N+n')c\le48x^{49}E_1\le x^{51}E_1.
\]
The left side of (GI.9) is at most $x^{86}E_1E_2$.
The final four inequalities of (GI.5) give
\[
 85\log x+\tfrac12x^{\varepsilon/4}
       +4x^{23\varepsilon/48}+x^{\varepsilon/3}
       \le x^\varepsilon.
\]
Exponentiating proves $x^{86}E_1E_2\le
xe^{x^\varepsilon-x^{\varepsilon/3}}=b'$ and completes
the literal fourteen-part hypothesis.

\paragraph{GI5: the actual family statement.}
For a supplied \texttt{FamP}, the window hypothesis is the
eventual inequality $N_1\le K.\mathrm{CN}\log x$.
Substitution into the nonnegative second summand of $\rho$ gives
$\rho\le\rho_B\le x^{-2}$, where the last bound is exactly
\texttt{FamP.rho\_le}. This uses no prime sum over the excluded
band. The family supplies (GI.3), and its filling identity
$\lambda=L.\mathrm{LaG}$ makes the conductor inequality and
the chosen lower-degree window exact.

Indeed, at every selected position the actual integer degrees
lie in
\[
 W_a=[\lceil a\rceil,\lfloor a+a/(kY)\rfloor]\cap\mathbb Z,
 \qquad
 W_b=[\lceil b'\rceil,\lfloor b'+I\rfloor]\cap\mathbb Z.  \tag{GI.10}
\]
This follows from \texttt{aT\_win}, \texttt{bT\_ge} and
\texttt{SizeBot}, followed by the definition of the orbit
restriction. Combining these memberships with (GI.6)--(GI.9)
proves the exact \texttt{Joints.HypIAt}: eventually, for
\emph{every} family index, candidate Type-I prime and good
representative, the two spacing bounds, numerical hypothesis
and actual weight windows hold. A single maximum of the fixed
thresholds suffices; no choice depending on the prime or index
has been made. This does not assert existence of a family.

\paragraph{GI6: convert to the completed local theorem.}
Take
\[
\begin{gathered}
 C_0=K.\mathrm{La},\quad
 0\le C_H\le\min\{K.\mathrm{kI},K.\mathrm{DI}\},\quad
 2C_H\le K.\rho,\\
 m=k-1,\quad c_{\rm cl}=\lfloor c\rfloor_+,\quad
 D_p=D_f,\quad c_*=8(D_f+1),\quad
 b_l=b',\quad b_u=b'+I.
\end{gathered}                                                    \tag{GI.11}
\]
GII1 gives an even selected point, at least two segments, and
their bounds $\mathrm{segG}\le G'$, $\mathrm{segS}\le N_1$.
The only delicate normalization is the local band radius:
\[
 \delta_I=\frac{C_HN_1\lambda}k+\frac{2I}{b'},\qquad
 2\delta_I\le\rho<1/6.                                  \tag{GI.12}
\]
Hence the local forbidden interval
$[X'(1-2\delta_I),X'(1+2\delta_I)]$ is contained in the
global interval with radius $\rho$. Exclusion of
\texttt{InBand0} therefore excludes every odd prime power
required locally. An actual candidate excludes the larger
$\rho_B$ interval, so supplies this premise.

All remaining local fields follow from (GI.7)--(GI.9):
$m+1=k$, $N_1\ge1$, $I\le b'/2$, $c_{\rm cl}\le c$,
and
\[
 C_HNY^2(N_1^5+k^3)\log_p a\le D_f
\]
by $N\le N_+$ and $C_H\le K.\mathrm{DI}$.
The last size product only decreases on replacing
$K.\mathrm{kI},c$ by $C_H,c_{\rm cl}$. This proves the
literal local \texttt{HypI}, with its original depth and
uniformizer normalization.

\paragraph{GI7: compare the exact time bounds.}
For each odd $y\in T$, its nearest preceding even selected
point has a unique least backward distance $1\le d<n'$.
The forward and backward descriptions are equivalent modulo
$n'$. Consequently the local \texttt{dPrev,yPrev} and global
\texttt{dMinus} define the same leading sum
\[
 S=\sum_{y\in T,\ y\text{ odd}}
        \min\{b_y,a_{y-d}/p^d\}.
\]
This identity holds for arbitrary signed integer degrees.
Put $E=a/p^2+\lambda a/(pk)$ and
\[
\begin{aligned}
 c_{\rm loc}&=8(D_f+1)+CNk\log_p a+4Nc_{\rm cl}+2n'c_{\rm cl},\\
 c_6&=8(D_f+1)+K.\kappa_2N_+k\log_p a+4N_+c+2n'c,
\end{aligned}
\]
where $K.\kappa_1,K.\kappa_2$ abbreviate the two global
\texttt{kappaThreshold} coefficients. The actual definitions are
\[
\begin{aligned}
 t_I={}&S+CNN_1(E+b'/p+\lambda b'/k)\\
       &+CN(Qc_{\rm loc}+kQ+k\log_p a+c_{\rm cl})+2n'c_{\rm cl},\\
 t_{If}={}&S+K.\kappa_1NN_1E\\
       &+K.\kappa_2N(Qc_6+kQ+k\log_p a+c)+2n'c.
\end{aligned}                                                    \tag{GI.13}
\]
If $2C\le K.\kappa_1$ and $C\le K.\kappa_2$, then
$b'\le a/p$, $c_{\rm loc}\le c_6$ and $c_{\rm cl}\le c$
prove term by term
$t_I\le t_{If}\le\lceil t_{If}\rceil_+$.
The common signed sum $S$ is unchanged.

\paragraph{GI8: the actual exceptional mean.}
TC6 supplies one function $\kappa$ on $W_a^T\times W_b^T$,
chosen before the coefficient rings, with mean at most
$CNN_1/Y$. GI7 makes all its certificates valid at the exact
global time $\lceil t_{If}\rceil_+$. Leastness of the source's
canonical exceptional function therefore gives
\[
 0\le\kappa_I(\mathbf a,\mathbf b)
       \le\kappa(\mathbf a|_T,\mathbf b|_T).
\]
Restriction to $T$ is a bijection from the source's orbit box
to $W^T$; its inverse extends by zero. Thus both double sums
and their denominators agree, proving the literal
\texttt{adm8} bound
\[
 \operatorname{meanLoc}(\kappa_I)\le C\,NN_1/Y.           \tag{GI.14}
\]
GJ1 supplies attainment and unpacks the actual family field
\texttt{ex.kapI\_def}; an arbitrary instance must still satisfy
\texttt{KapIDef}. Empty boxes give zero under the source convention.

The following fixed vector is sufficient for TC6 and every
comparison above:
\[
\begin{gathered}
 K.\mathrm{La}\ge880,\qquad
 K.\mathrm{kI},K.\mathrm{DI}\ge2800001,\qquad
 K.\rho\ge5600002,\\
 K.\kappa_1\ge604228,\qquad K.\kappa_2\ge302114,
 \qquad C=302114.
\end{gathered}                                                    \tag{GI.15}
\]
The mean coefficient may be any $C_m\ge302114$.
These choices precede the scale, arithmetic object, prime,
weights and coefficient rings. Global prime averaging and the
existence of the simultaneous arithmetic family remain separate.
